%!TEX program = pdflatex
% 通用学术论文LaTeX模板
% 适用于大多数会议和期刊

\documentclass[12pt,a4paper]{article}

% 编码设置
\usepackage[utf8]{inputenc}
\usepackage[T1]{fontenc}

% 页面设置
\usepackage[margin=1in]{geometry}
\usepackage{setspace}
\setstretch{1.5}

% 数学支持
\usepackage{amsmath}
\usepackage{amssymb}
\usepackage{amsfonts}

% 图表支持
\usepackage{graphicx}
\usepackage{subcaption}
\usepackage{float}
\usepackage{booktabs}

% 表格增强
\usepackage{tabularx}
\usepackage{longtable}
\usepackage{colortbl}

% 代码环境
\usepackage{listings}
\usepackage{verbatim}

% 参考文献
\usepackage{natbib}
\bibliographystyle{plainnat}

% 超链接
\usepackage[colorlinks=true, linkcolor=blue, citecolor=blue, urlcolor=blue]{hyperref}

% 定理环境
\usepackage{amsthm}
\newtheorem{theorem}{Theorem}
\newtheorem lemma{Lemma}
\newtheorem definition{Definition}
\newtheorem proposition{Proposition}
\newtheorem corollary{Corollary}
\newtheorem proof{Proof}

% 页眉页脚
\usepackage{fancyhdr}
\pagestyle{fancy}
\fancyhf{}
\rhead{\thepage}
\lhead{\leftmark}

% 自定义命令
\newcommand{\R}{\mathbb{R}}
\newcommand{\N}{\mathbb{N}}
\newcommand{\Z}{\mathbb{Z}}
\newcommand{\E}{\mathbb{E}}
\DeclareMathOperator*{\argmax}{argmax}
\DeclareMathOperator*{\argmin}{argmin}

% Title
\title{Your Paper Title Here}
\author{Author Name$^1$, Co-Author Name$^2$ \\
$^1$Department of Computer Science, University Name \\
$^2$Another Institution \\
\{author1, author2\}@example.com}

\begin{document}

\maketitle

\begin{abstract}
Your abstract goes here. It should summarize the main contributions of your paper in a concise manner. Keep it between 150-300 words for a typical conference paper.
\end{abstract}

\section{Introduction}
\label{sec:introduction}

Introduce the problem and motivation here. Explain why this work is important and what gap it fills in the literature.

\subsection{Contributions}

Summarize your main contributions:

\begin{itemize}
    \item First contribution
    \item Second contribution
    \item Third contribution
\end{itemize}

\section{Related Work}
\label{sec:related}

Discuss related work here. Compare and contrast your approach with existing methods.

\section{Methodology}
\label{sec:method}

Explain your methodology in detail.

\subsection{Background}

Provide necessary background information.

\subsection{Proposed Approach}

Describe your proposed approach in detail.

\subsection{Theoretical Analysis}

Provide theoretical analysis if applicable.

\section{Experiments}
\label{sec:experiments}

Describe your experimental setup and results.

\subsection{Datasets}

Describe the datasets used in your experiments.

\subsection{Baselines}

List the baseline methods you compare against.

\subsection{Results}

Present your experimental results.

\begin{table}[htbp]
    \centering
    \caption{Example table}
    \label{tab:example}
    \begin{tabular}{lcc}
        \toprule
        Method & Accuracy & F1-Score \\
        \midrule
        Baseline 1 & 85.2 & 0.85 \\
        Baseline 2 & 87.1 & 0.87 \\
        Ours & \textbf{89.5} & \textbf{0.89} \\
        \bottomrule
    \end{tabular}
\end{table}

\begin{figure}[htbp]
    \centering
    \includegraphics[width=0.8\textwidth]{figures/fig1.png}
    \caption{Example figure}
    \label{fig:example}
\end{figure}

\section{Conclusion}
\label{sec:conclusion}

Summarize your work and discuss future directions.

\section{Acknowledgments}
\label{sec:acks}

Thank collaborators and funding agencies.

% References
\bibliography{references}

\end{document}
